\section{Arquitectura del sistema y distribución de microcontroladores}

La distribución actual del sistema utiliza tres ATmega328P conectados mediante un bus I2C compartido. El Maestro coordina la red, controla la pantalla LCD y se comunica directamente con el VL53L0X. Los dos Slaves tienen direcciones I2C distintas y concentran sensores y actuadores de diferentes zonas del proceso.

\subsection{ATmega328P Maestro}
\begin{itemize}[leftmargin=*]
    \item UART por PD0/RX y PD1/TX, configurada actualmente a 9600 baudios, 8N1.
    \item LCD de 16x2 en modo de 8 bits: D0--D7 conectados desde PD2 hasta PB1; E en PC0 y RS en PC1.
    \item Bus I2C: SDA en PC4/A4 y SCL en PC5/A5.
    \item VL53L0X (Slave 3) en el bus I2C, con dirección 0x29.
    \item XSHUT del VL53L0X conectado a PB2/D10.
\end{itemize}

\subsection{ATmega328P Slave 1 -- dirección I2C 0x11}
\begin{itemize}[leftmargin=*]
    \item HC-SR04: ECHO en PC2/A2 y TRIG en PC3/A3.
    \item Servo de agua en PB3/D11, con posiciones actuales de 0 grados para cerrado y 90 grados para abierto.
    \item L298N para el motor DC: ENA en PB1/D9, IN1 en PD7/D7 e IN2 en PB0/D8.
    \item PWM de lavado indicado actualmente: 128.
\end{itemize}

\subsection{ATmega328P Slave 2 -- dirección I2C 0x12}
\begin{itemize}[leftmargin=*]
    \item Sensor infrarrojo en PC0/A0, activo en nivel bajo.
    \item ULN2003 del Stepper 28BYJ-48: IN1--IN4 conectados a PD2/D2, PD3/D3, PD4/D4 y PD5/D5.
    \item Servo de puerta en PD6/D6, con posiciones actuales de 145 grados para cerrado y 50 grados para abierto.
    \item Movimiento del Stepper mediante medio paso, con velocidad indicada actualmente de 800 pasos por segundo.
\end{itemize}

\subsection{Bus I2C compartido}
SDA conecta A4 del Maestro, A4 del Slave 1, A4 del Slave 2 y SDA del VL53L0X. SCL conecta de forma equivalente los pines A5 y SCL de los dispositivos. Se conectan resistencias de 470~\(\Omega\) en configuración \textit{pull-up} en las líneas SDA y SCL.

\subsection{Alimentación}
Los microcontroladores, sensores y señales de control comparten una referencia común de tierra. El motor DC utiliza una fuente independiente conectada a VS del L298N; los servomotores deben alimentarse mediante una fuente externa regulada de 5 V con capacidad de corriente suficiente; y el Stepper debe alimentarse a través del ULN2003 con una fuente adecuada. Todas las tierras deben compartir referencia, procurando evitar que las corrientes de los motores circulen por el cableado del bus I2C.
